\documentclass[10pt]{article}
\usepackage[T1]{fontenc}
\usepackage[utf8]{inputenc}
\usepackage[french]{babel}
\usepackage{amsmath,amssymb,mathrsfs,mathpazo}
\usepackage{geometry}
\geometry{a4paper,margin=2cm}
\usepackage{graphicx,array,tabularx,tabularray,colortbl,multirow,multicol,booktabs,enumitem,pifont,ulem,fancybox,diagbox,cancel}
\usepackage{pstricks,pstricks-add,pst-plot,pst-tree,pst-node,pst-text,pst-func,pst-math,pst-bspline,pst-3dplot,pst-eucl,pst-solides3d}
\let\pscalculate\undefined
\usepackage{pst-bezier}
\usepackage{tikz,pgfplots,tkz-tab}
\usetikzlibrary{arrows,arrows.meta,patterns,calc,positioning,trees,decorations.markings,shapes.geometric,shapes.misc,shapes.arrows,intersections,decorations.pathmorphing,decorations.pathreplacing,backgrounds,shadows,babel}
\pgfplotsset{compat=1.18}
\usepgfplotslibrary{fillbetween}
\usepackage[np]{numprint}
\usepackage{listings,makecell,pas-tableur,twoopt}
\let\barre\undefined
\usepackage[pstricks]{tablvar}
\definecolor{darkgreen}{RGB}{0,100,0}
\providecommand{\red}{\color{red}}
\providecommand{\blue}{\color{blue}}
\providecommand{\green}{\color{green}}
\providecommand{\black}{\color{black}}
\providecommand{\white}{\color{white}}
\providecommand{\gray}{\color{gray}}
\providecommand{\cyan}{\color{cyan}}
\providecommand{\magenta}{\color{magenta}}
\providecommand{\yellow}{\color{yellow}}
\providecommand{\violet}{\color{violet}}
\providecommand{\orange}{\color{orange}}
\providecommand{\euro}{\texteuro}
\providecommand{\EUR}{\texteuro}
\providecommand{\ieme}{\textsuperscript{e}}
\providecommand{\degres}{\ensuremath{^\circ}}
\providecommand{\degree}{\ensuremath{^\circ}}
\providecommand{\vv}[1]{\overrightarrow{#1}}
\providecommand{\Coord}[1]{#1}
\providecommand{\decosix}{}
\providecommand{\starredbullet}{\ensuremath{\bullet}}
\pagestyle{empty}
\setlength{\parindent}{0pt}
\usepackage{adjustbox}\DeclareUnicodeCharacter{2236}{:}\DeclareUnicodeCharacter{2190}{\ensuremath{\leftarrow}}
\providecommand{\R}{\mathbb{R}}
\providecommand{\N}{\mathbb{N}}
\providecommand{\Z}{\mathbb{Z}}
\providecommand{\Q}{\mathbb{Q}}
\providecommand{\D}{\mathbb{D}}
\providecommand{\C}{\mathbb{C}}
\providecommand{\e}{{\mathrm{e}}}
\providecommand{\vect}[1]{\overrightarrow{#1}}
\providecommand{\vectt}[1]{\overrightarrow{\mathrm{#1}}}
\providecommand{\vvt}[1]{\overrightarrow{\mathrm{#1}}}
\providecommand{\np}[2][]{{#2~#1}}
\providecommand{\vv}[1]{{\overrightarrow{#1}}}
\renewcommand{\d}{{\mathrm{d}}}
\providecommand{\degres}{{^\circ}}
\providecommand{\degre}{{^\circ}}
\providecommand{\blue}{\color{blue}}
\providecommand{\red}{\color{red}}
\providecommand{\enskip}{\quad}
\providecommand{\llbracket}{\lbrack\!\lbrack}
\providecommand{\rrbracket}{\rbrack\!\rbrack}
\providecommand{\euro}{{\text{€}}}
\providecommand{\no}{{\text{n°}}}
\providecommand{\interval}[2]{{[#1;#2]}}
\providecommand{\intervFO}[2]{{[#1;#2[}}
\providecommand{\dotfill}{\ldots}

\providecommand{\ds}{\displaystyle}
\providecommand{\textding}[1]{#1}
\providecommand{\Oij}{\(\mathrm{O};\vec{\imath},\vec{\jmath}\)}
\providecommand{\Oijk}{\(\mathrm{O};\vec{\imath},\vec{\jmath},\vec{k}\)}
\providecommand{\Ouv}{\(\mathrm{O};\vec{u},\vec{v}\)}
\providecommand{\intervFF}[2]{\ensuremath{[#1;#2]}}
\renewcommand{\euro}{\ifmmode\text{\texteuro}\else\texteuro\fi}
\DeclareUnicodeCharacter{00B0}{\ifmmode^{\circ}\else\textdegree\fi}
\DeclareUnicodeCharacter{20AC}{\ifmmode\text{\texteuro}\else\texteuro\fi}
% Énoncés : Éducation nationale. Transcription et corrigés : APMEP, Denis Vergès et contributeurs. Découpage et classement : CoopMaths, Sébastien Lozano, Éric Elter et contributeurs. Adaptation de la présentation : Maths Explorer.
% CC BY-SA : https://coopmaths.fr/www/about/
% https://coopmaths.fr/alea/static/eam/2026/tex/eam_technologique_2026_06_polynesie_2.tex
\begin{document}
\noindent\textit{APMEP / CoopMaths. CC BY-SA. Présentation adaptée par Maths Explorer.}\par\medskip
Une personne souhaite emprunter \np{10000}~\euro{}, remboursables en six versements annuels.

Un établissement de crédit propose les deux formules.

\begin{itemize}[label={\textbullet}]
\item \textbf{Formule 1 :} Le premier versement est \np{1025}~\euro{} et les versements annuels suivants augmentent de \np{400}~\euro{} chaque année par rapport au versement précédent.
\item \textbf{Formule 2 :} le premier versement est \np{1550}~\euro{} et les versements annuels suivants augmentent de 10~\% chaque année par rapport au versement précédent.
\end{itemize}

\subsection*{Partie A}

On se place dans le cas de la \textbf{formule 1}.

Étant donné un entier naturel $n$, on appelle $u_n$ le versement, exprimé en euros, de l'année $2025+n$.

On a \quad $u_0=\np{1025}$.

\begin{enumerate}
\item Calculer $u_1$.
\item Donner la nature de la suite $(u_n)$ et donner sa raison.
\item Donner la valeur du dernier versement $u_5$.
\end{enumerate}

\subsubsection*{Partie B}

On se place dans le cas de la \textbf{formule 2}.

Étant donné un entier naturel $n$, on appelle $v_n$ le versement, exprimé en euros, de l'année $2025+n$.

On a \quad $v_0=\np{1550}$.

\begin{enumerate}
\item Quel calcul permet d'obtenir $v_1=\np{1705}$ ?
\item Quelle est la nature de la suite $(v_n)$ et donner sa raison.
\end{enumerate}

\subsubsection*{Partie C}

La feuille de calcul suivante présente, pour chaque année et pour chaque formule, le cumul des montants remboursés depuis le début du crédit.

Le montant final en gras dans le tableau, en ligne 7, correspond à la somme des \np{10000}~\euro{} empruntés et des intérêts perçus par l'établissement de crédit.

\begin{center}
\helvbx
\begin{tblr}{
	colspec={c},
	row{1} = {gray9,c},
	column{1} = {gray9,c,5mm},
	column{2} = {c,15mm},
	column{3} = {c,25mm},
	column{4} = {c,25mm},
	hlines,vlines,
	}
 	& 	A 	& 	B 								& 	C 								\\
1	&	$n$	&	\textbf{Formule 1}				&	\textbf{Formule 2} 				\\
2	&	0	&	\np{1025.00}~\euro{}				&	\np{1550.00}~\euro{}				\\
3	&	1	&	\np{2450.00}~\euro{}				&	\np{3255.00}~\euro{}				\\
4	&	2	&	\np{4275.00}~\euro{}				&	\np{5130.50}~\euro{}				\\
5	&	3	&	\np{6500.00}~\euro{}				&	\np{7193.55}~\euro{}				\\
6	&	4	&	\np{9125.00}~\euro{}				&	\np{9462.91}~\euro{}				\\
7	&	5	&	\textbf{\np{12150.00}~\euro{}}	&	\textbf{\np{11959.20}~\euro{}}	\\
\end{tblr}
\end{center}

Quelle est la formule la plus avantageuse pour la personne qui emprunte ?

Justifier.

\end{document}
